%% 12-key-performance-indicators-and-value-indicators.tex -- generated by tools/build_from_docx.py; edit freely

\hypertarget{key-performance-indicators-and-value-indicators}{%
\section{Key Performance Indicators and Value Indicators}\label{key-performance-indicators-and-value-indicators}}

The preceding chapters each state what their subject contributes to the
project's indicators. This chapter does something those tables cannot:
it accounts for what the framework costs.

The premise is that the framework does not merely host algorithms that
meet project targets. If the platform itself fails on latency,
scalability or reliability, the targets are missed regardless of how
good the contributed algorithms are. Platform overhead is therefore a
first-class design constraint rather than an implementation detail, and
this chapter states it as one.

Three things follow. The indicators that directly constrain the
framework are identified, with what binds about each. Every guardrail
specified in the preceding chapters is classified by where its cost
falls. And the two places where a guardrail and an indicator are in
genuine tension are resolved explicitly, rather than left for an
implementer to discover.

\hypertarget{the-primary-commitment}{%
\subsection{The Primary Commitment}\label{the-primary-commitment}}

The project defines the reduction of real-time control-loop timescales
below ten milliseconds as its own indicator, and its definition of a
control loop is decisive for this framework: the elapsed time across
observation and measurement collection from the radio access network,
decision making by the control function, and enforcement of the control
action on the network.

That is exactly the path this framework specifies. The baseline is the
ten millisecond to one second range of the near-real-time controller in
the reference architecture; sub-ten-millisecond operation is a target on
top of it.

Consequently the cost classification of Section 12.3 is not an internal
preference. It is the mechanism by which the commitment is met:
inference on the per-decision path, and every other guardrail off it.

\hypertarget{the-indicators-that-bind}{%
\subsection{The Indicators that Bind}\label{the-indicators-that-bind}}

Several indicators constrain this framework directly, and each binds in
a different way.

Reconfiguration within ten milliseconds from the moment a decision is
taken is the binding constraint on the actuation path, and the wording
matters: the budget starts when the decision is taken and ends at
reconfiguration. The project's own stated mitigation is that decisions
be prepared in advance rather than triggered by observed failures,
staging actions proactively to relieve time pressure on the actuation
path. Section 12.4 specifies the mechanism.

Cross-domain end-to-end latency below twenty milliseconds bounds the
whole federated loop. Cross-domain jitter within two milliseconds, with
fewer than ten percent deadline misses, bounds variance rather than the
mean, which is what platform gates and retries actually threaten.
Reduced reaction time from identification to resolution charges loop
overhead directly.

Reduction of the conflict performance penalty by twenty percent is owned
by the conflict fabric. Service continuity at seven nines bounds
platform availability, which means the fabric, the registry and the
gates must not become the least available element in the system. The
requirement that twin-validated outputs not break time-sensitive
operation in ninety-nine percent of cases ratifies guarded deployment
directly, since it calls for unsafe or low-confidence actions to be
rejected or replaced by a baseline.

Service creation and termination time, link utilization and operational
expenditure are bounded by lifecycle and admission cost. The energy
indicators matter here too, because platform compute is part of the
energy budget, and a guardrail that runs every cycle spends energy every
cycle. Table 12-1 states the indicators that bind and what each
constrains.

\begin{wptable}{12-1}

\begin{longtable}[]{@{}
  >{\raggedright\arraybackslash}p{(\columnwidth - 4\tabcolsep) * \real{0.1750}}
  >{\raggedright\arraybackslash}p{(\columnwidth - 4\tabcolsep) * \real{0.3375}}
  >{\raggedright\arraybackslash}p{(\columnwidth - 4\tabcolsep) * \real{0.4875}}@{}}
\caption{Indicators that constrain the framework
directly}\tabularnewline
\toprule\noalign{}
\begin{minipage}[b]{\linewidth}\raggedright
\textbf{KPI}
\end{minipage} & \begin{minipage}[b]{\linewidth}\raggedright
\textbf{Target}
\end{minipage} & \begin{minipage}[b]{\linewidth}\raggedright
\textbf{What it binds}
\end{minipage} \\
\midrule\noalign{}
\endfirsthead
\toprule\noalign{}
\begin{minipage}[b]{\linewidth}\raggedright
\textbf{KPI}
\end{minipage} & \begin{minipage}[b]{\linewidth}\raggedright
\textbf{Target}
\end{minipage} & \begin{minipage}[b]{\linewidth}\raggedright
\textbf{What it binds}
\end{minipage} \\
\midrule\noalign{}
\endhead
\bottomrule\noalign{}
\endlastfoot
\textbf{KPI\#34} & Real-time control-loop timescale below 10 ms. & The
primary commitment. Its definition of a loop is the MS, AE, DE, ACT path
this framework specifies. \\
\textbf{KPI\#4} & Domain reconfiguration within 10 ms from the moment a
decision is taken. & Binds the actuation path. Budget starts at the
decision, not at the observation. \\
\textbf{KPI\#7} & Cross-domain end-to-end latency below 20 ms. & Bounds
the whole federated loop. \\
\textbf{KPI\#8} & Cross-domain jitter within 2 ms, deadline misses below
10 percent. & Bounds variance, which gates and retries threaten more
than they threaten the mean. \\
\textbf{KPI\#19} & Reduced reaction time from identification to
resolution. & Charges loop overhead directly. \\
\textbf{KPI\#21} & Conflict performance penalty reduced by 20 percent. &
Owned by the conflict fabric; measured from the verdict counters of
Chapter 10. \\
\textbf{KPI\#37} & Service continuity at 99.99999 percent. & The fabric,
registry and gates must not be the least available element. \\
\textbf{KPI\#42} & Twin-validated outputs must not break time-sensitive
operation in 99 percent of cases. & Ratifies guarded deployment: unsafe
or low-confidence actions rejected or replaced by a baseline. \\
\textbf{KPI\#6, KPI\#20, KPI\#38} & Service creation and termination
time, link utilization, operational expenditure. & Bounded by lifecycle
and admission cost. \\
\textbf{KPI\#43, KPI\#16, KPI\#14} & Energy efficiency, radio link
power, energy cost. & Platform compute is part of the energy budget; a
per-cycle guardrail spends energy per cycle. \\
\end{longtable}

\end{wptable}

\hypertarget{guardrail-cost-classes}{%
\subsection{Guardrail Cost Classes}\label{guardrail-cost-classes}}

Every guardrail in this framework is classified by where its cost falls.
The classification is normative: an implementation may not relocate a
guardrail into a more expensive class without recomputing the effective
latency of the affected agents.

Class A is admission-time only, with zero per-cycle cost. Schema,
compatibility, integrity and attestation checks; scope nesting and
delegation coverage; enumerated partial capabilities; the contract-shape
prohibitions; machine-evaluability of validity predicates; typed failure
policy; conflict policy and priority registration; action-space
bounding; and control-parameter registration with the fabric. These are
the bulk of the framework's checking, deliberately placed here.

Class B is per-message, local and constant-time, of sub-microsecond
order. Envelope field population; bus-boundary type and schema
validation; provenance stamping; generation fencing; degraded-evidence
gate evaluation; last-mile bounds re-validation; retry counters and rate
limits; and feedback emission. These are arithmetic and field checks
over data already in hand, with no network hop and no blocking call.
They are charged to the loop budget as fixed local overhead, and they
are not tiered away for fast loops, because their cost is not what
threatens the latency targets.

Class C is per-decision and involves a network hop, of milliseconds
order. Conflict mediation is the only member. It is the only guardrail
placing a remote synchronous call in the decision path, and therefore
the only one subject to the envelope tiering of Section 12.4.

Class D is per-actuation and involves the southbound path, with variable
cost. Read-back verification and rollback. These are subject to the
verification tiering of Section 12.5. Table 12-2 states the classes with
their members and per-cycle cost.

\begin{wptable}{12-2}

\begin{longtable}[]{@{}
  >{\raggedright\arraybackslash}p{(\columnwidth - 6\tabcolsep) * \real{0.0562}}
  >{\raggedright\arraybackslash}p{(\columnwidth - 6\tabcolsep) * \real{0.2125}}
  >{\raggedright\arraybackslash}p{(\columnwidth - 6\tabcolsep) * \real{0.5187}}
  >{\raggedright\arraybackslash}p{(\columnwidth - 6\tabcolsep) * \real{0.2126}}@{}}
\caption{Guardrail cost classes}\tabularnewline
\toprule\noalign{}
\begin{minipage}[b]{\linewidth}\raggedright
\textbf{Class}
\end{minipage} & \begin{minipage}[b]{\linewidth}\raggedright
\textbf{Where the cost falls}
\end{minipage} & \begin{minipage}[b]{\linewidth}\raggedright
\textbf{Members}
\end{minipage} & \begin{minipage}[b]{\linewidth}\raggedright
\textbf{Per-cycle cost}
\end{minipage} \\
\midrule\noalign{}
\endfirsthead
\toprule\noalign{}
\begin{minipage}[b]{\linewidth}\raggedright
\textbf{Class}
\end{minipage} & \begin{minipage}[b]{\linewidth}\raggedright
\textbf{Where the cost falls}
\end{minipage} & \begin{minipage}[b]{\linewidth}\raggedright
\textbf{Members}
\end{minipage} & \begin{minipage}[b]{\linewidth}\raggedright
\textbf{Per-cycle cost}
\end{minipage} \\
\midrule\noalign{}
\endhead
\bottomrule\noalign{}
\endlastfoot
\textbf{A} & Admission time only & Schema, compatibility, integrity and
attestation; scope nesting and delegation coverage; enumerated partial
capabilities; contract-shape prohibitions; machine-evaluable validity
predicates; typed failure policy; conflict and priority registration;
action-space bounding; CP registration. & Zero per cycle \\
\textbf{B} & Per message, local, constant time & Envelope field
population; bus-boundary type and schema validation; provenance
stamping; generation fencing; degraded-evidence gate; last-mile bounds
re-validation; retry counters and rate limits; feedback emission. &
Sub-microsecond order; never tiered away \\
\textbf{C} & Per decision, network hop & Conflict mediation. &
Milliseconds order; subject to envelope tiering \\
\textbf{D} & Per actuation, southbound path & Read-back verification;
rollback. & Variable; subject to verification tiering \\
\end{longtable}

\end{wptable}

The distinction prevents two opposite errors: weakening cheap local
checks for an imagined performance benefit, and assuming expensive
remote checks are as free as local ones. It also tells an implementer
precisely where optimization effort is worth spending. The classes map
onto the terms of the effective loop latency of Section 3.8.3: Class B
is the platform overhead term, Class C the mediation term, zero for an
action type on the envelope or policy path, and Class D the verification
term, zero where verification is deferred. Class A contributes nothing
to that sum, which is why it carries the bulk of the checking.

An agent whose declared budget cannot absorb its own guardrails is
rejected, or is assigned a lower effective class, as its class policy
declares (Section 3.8.3). It is never admitted with the guardrails
silently dropped. A safety mechanism whose cost is unbudgeted does not
make the system safe; it makes the system late, and lateness is itself a
failure against these indicators.

\hypertarget{envelope-tiering-resolves-the-mediation-tension}{%
\subsection{Envelope Tiering Resolves the Mediation Tension}\label{envelope-tiering-resolves-the-mediation-tension}}

Synchronous fail-closed mediation is in direct tension with a
sub-ten-millisecond actuation budget. The resolution separates the
authorizing decision from the executing action, and it weakens no
invariant.

For hard-real-time loops the fabric issues a verdict over an action
envelope, meaning a bounded region of the action space with a validity
window and a rate limit, rather than over each individual command. The
Decision Engine then actuates within that pre-authorized envelope at
line rate, with no per-command mediation hop. Actuation outside the
envelope, or after its validity window expires, requires a fresh
synchronous verdict.

The envelope is itself a verdict. Authority still originates with the
fabric and never with the agent, so the prohibition on self-approval
holds exactly as stated: an agent cannot authorize itself, and can only
act inside authority already granted.

This is the project's own stated mitigation for the reconfiguration
target, and it is how policy guidance already works in the reference
radio architecture, where a non-real-time controller issues bounding
policy on a slow loop and a near-real-time controller acts quickly
within it. Envelope grants, their bounds and their expiry are held in
the registry, and last-mile re-validation checks each command against
the live envelope rather than a cached copy. This is the action envelope
of Section 3.9.2; the distributed policy of the same section gives the
same shape where the permitted region cannot be enumerated in advance.

Per-command mediation on a ten millisecond budget would either miss the
target or be bypassed in practice, and the second outcome is worse than
the first. Tiering preserves the invariant at the timescale where it is
enforceable and preserves the indicator at the timescale where that
matters.

\hypertarget{verification-tiering-protects-the-jitter-bound}{%
\subsection{Verification Tiering Protects the Jitter Bound}\label{verification-tiering-protects-the-jitter-bound}}

Mandatory read-back must not serialize inside a jitter-bounded loop.

For each action type, the action space declares whether verification is
in band or deferred and, where deferred, the deadline within which it
completes (Section 3.5.8). In band, the transaction is not closed until
the change is confirmed, which suits slow loops. Deferred, the command
proceeds and verification completes within the declared deadline, which
is what a hard-real-time loop requires. The action reports
acknowledged\_\allowbreak{}unverified until confirmed, and a missed
deadline is a conformance fault that feeds Decision Engine re-decision
and assurance.

This is the same pattern as the action envelope rather than a second
mechanism. The rule governing verification is granted in advance, with
the admitted descriptor, and held by the Actuator's runtime, which
enforces the deadline locally without a round trip; nothing watches from
outside. What distinguishes the tiers is therefore not synchrony against
asynchrony, but whether confirmation is part of the transaction or bound
to it by a deadline granted in advance.

The invariant is preserved. An acknowledgement is still never proof of
change. Tiering changes when the system learns, never whether it checks.

Blocking a ten millisecond loop on a read-back round trip would breach
the reconfiguration target and inflate jitter against the variance
bound. Deferring the check costs a bounded window of uncertainty, which
is declared, counted and visible in assurance. That is an honest
measured exposure in place of a hidden one, which is the trade this
framework makes repeatedly.

\hypertarget{scalability-is-bounded-by-design}{%
\subsection{Scalability is Bounded by Design}\label{scalability-is-bounded-by-design}}

Seven nines of continuity and a cross-domain latency bound together make
platform components a potential single point of failure and a potential
bottleneck. Four properties bound this by design.

Mediation is per-domain rather than global, which is what keeps conflict
resolution horizontally scalable as domains multiply. Per-service and
per-agent topics scale with services and agents rather than with the
square of the agent count, and because agents do not cross-subscribe,
message load grows linearly. Registry reads on the intent-egress path
are served from a replica invalidated on every change of generation or
grant, so authority checks do not serialize on a central store and never
act on a superseded grant. And admission cost is paid once at
instantiation rather than per cycle.

That last property is the reason so many closures in this framework were
placed at admission rather than at runtime. It is what protects service
creation time from guardrail overhead, and it is stated here so that
future partners do not relocate checks into the hot path without seeing
what they would break.

\hypertarget{degradation-must-be-cheaper-than-failure}{%
\subsection{Degradation Must Be Cheaper than Failure}\label{degradation-must-be-cheaper-than-failure}}

Fail-closed behavior protects correctness but can threaten continuity if
applied bluntly. Suspending every agent on a fabric hiccup would itself
breach the continuity target.

Loss of the mediation binding therefore suspends decision emission, not
the service: existing configuration persists and the data plane keeps
running. The agent moves to suspended rather than terminating, with its
monitoring and analytics continuing, so recovery restores an instance
rather than re-instantiating it, which would cost both service creation
time and recovery time. An action envelope granted before the loss is
not voided by the fabric's unavailability and remains valid until its
declared expiry (Section 3.9.2), so an agent holding one continues to
act within it and is suspended at that expiry if the binding has not
been restored; an agent holding none is suspended at once. A brief
fabric outage therefore does not immediately stop a hard-real-time loop.

Fail closed must mean stop deciding, never stop carrying traffic.
Conflating the two would convert a control-plane fault into a service
outage, trading a correctness guarantee for a continuity breach, which
is a bad trade in every operational circumstance this framework
anticipates.

\hypertarget{value-indicators-are-addressed-structurally}{%
\subsection{Value Indicators are Addressed Structurally}\label{value-indicators-are-addressed-structurally}}

The project names safety, security, trustworthiness, inclusiveness and
sustainability as value indicators. Each is discharged by mechanisms
already specified rather than by assertion.

Trustworthiness is platform-computed evidence, mandatory provenance,
verified actuation outcomes, registry-held authority and tamper-evident
conflict records. It is a set of checkable artifacts rather than a
claim. Safety is bounded action spaces, last-mile re-validation, the
prohibition on self-approval, and guarded deployment matching the
reject-or-baseline rule. Sustainability is platform overhead counted
against the energy indicators, which is the reason for front-loading
checks to admission rather than repeating them per cycle. Security is
trust-domain scoping and pseudonymization at egress, with signed
descriptors and artifact attestation at admission. Inclusiveness is the
platform premise itself: any partner may contribute an agent through the
same declared interface, including through the hardened extension route,
without privileged access. Table 12-3 states each value indicator
against the mechanisms that discharge it.

\begin{wptable}{12-3}

\begin{longtable}[]{@{}
  >{\raggedright\arraybackslash}p{(\columnwidth - 4\tabcolsep) * \real{0.1750}}
  >{\raggedright\arraybackslash}p{(\columnwidth - 4\tabcolsep) * \real{0.5125}}
  >{\raggedright\arraybackslash}p{(\columnwidth - 4\tabcolsep) * \real{0.3126}}@{}}
\caption{Value indicators and the mechanisms discharging
each}\tabularnewline
\toprule\noalign{}
\begin{minipage}[b]{\linewidth}\raggedright
\textbf{Value indicator}
\end{minipage} & \begin{minipage}[b]{\linewidth}\raggedright
\textbf{Mechanisms discharging it}
\end{minipage} & \begin{minipage}[b]{\linewidth}\raggedright
\textbf{Note}
\end{minipage} \\
\midrule\noalign{}
\endfirsthead
\toprule\noalign{}
\begin{minipage}[b]{\linewidth}\raggedright
\textbf{Value indicator}
\end{minipage} & \begin{minipage}[b]{\linewidth}\raggedright
\textbf{Mechanisms discharging it}
\end{minipage} & \begin{minipage}[b]{\linewidth}\raggedright
\textbf{Note}
\end{minipage} \\
\midrule\noalign{}
\endhead
\bottomrule\noalign{}
\endlastfoot
\textbf{Trustworthiness} & Platform-computed evidence, mandatory
provenance, verified actuation outcomes, registry-held authority,
tamper-evident conflict records. & A set of checkable artifacts rather
than a claim. \\
\textbf{Safety} & Bounded action spaces, last-mile re-validation, no
self-approval, guarded deployment on the reject-or-baseline rule. &
Discharges KPI\#42 directly. \\
\textbf{Sustainability} & Platform overhead counted against the energy
indicators; checks front-loaded to admission rather than repeated per
cycle. & Class A placement is the mechanism. \\
\textbf{Security} & Trust-domain scoping, pseudonymization at egress,
signed descriptors, artifact attestation at admission. & Enforced at
admission and at egress. \\
\textbf{Inclusiveness} & Any partner may contribute through the same
declared interface, including the hardened extension route, without
privileged access. & Multi-vendor participation as an outcome, not only
an architecture. \\
\end{longtable}

\end{wptable}

\hypertarget{two-value-indicators-justify-the-hardest-machinery}{%
\subsubsection{Two value indicators justify the hardest machinery}\label{two-value-indicators-justify-the-hardest-machinery}}

Two of the project's per-use-case value indicators bear directly on why
this framework's most expensive mechanisms earn their cost, and they are
stated by the project as human-value outcomes rather than engineering
niceties.

In the disaster use case, a named value indicator is the reduction of
miscoordination and threat to human life caused by service outages,
explicitly including outages caused by conflicting resource allocation
management processes, with citizens as the stakeholder. That is the
conflict fabric's reason to exist, expressed as a safety value rather
than as an architectural preference.

In the time-sensitive industrial use case, the named indicators are
reduced accidents at work due to communication failure, and reduced
energy per process unit.

So when a partner asks why verified actuation, fail-closed decision
making and conflict mediation justify their cost, the answer is in the
project's own value indicators. Outages from conflicting management
processes, and accidents from communication failure, are the harms these
guardrails exist to prevent.

\hypertarget{assurance-speaks-the-service-level-vocabulary}{%
\subsection{Assurance Speaks the Service-Level Vocabulary}\label{assurance-speaks-the-service-level-vocabulary}}

The project frames service level agreements as balancing deterministic,
stochastic, tail and average guarantees. The assurance this framework
returns to the requester does not collapse to averages.

A deadline-miss rate is a tail guarantee and a jitter bound is a
variance guarantee, and both are distinct from mean latency. The
four-outcome actuation vocabulary and deferred read-back verification
tiering preserve tail behavior rather than smoothing it away, which is
what makes the assurance report usable against a deterministic or tail
service level agreement rather than only against an average one.

\hypertarget{requirements-9}{%
\subsection{Requirements}\label{requirements-9}}

\begin{requirement}{Req-T23-KP-01}{MUST}{Every guardrail carries a declared cost class}
\reqfield{Description / rationale}{%
Every guardrail specified in this framework shall be classified as Class
A, B, C or D by where its cost falls. An implementation shall not
relocate a guardrail into a more expensive class without recomputing the
effective latency of every affected agent.}
\reqfield{Dependency}{%
Req-T23-FW-20}
\reqfield{Traceability (Forward)}{%
\begin{itemize}
\tightlist
\item
  Classification held with the guardrail specification.
\item
  Re-computation required on relocation.
\end{itemize}}
\reqfield{UNITY-6G components}{%
\begin{itemize}
\tightlist
\item
  Platform runtime
\item
  DMO
\end{itemize}}
\reqfield{Linked KPIs}{%
\begin{itemize}
\tightlist
\item
  KPI\#34: Reduce Open RAN real-time control loops time scales below 10
  ms and 10x in NTN for real-time 6G.
\item
  KPI\#4: Real-time reconfiguration of IEEE 802.11 network domain in
  range smaller than 10 ms from taking the decision.
\end{itemize}}
\end{requirement}

\begin{requirement}{Req-T23-KP-02}{MUST}{Effective latency covers the agent\textquotesingle s own guardrails}
\reqfield{Description / rationale}{%
The effective loop latency shall be computed as in Section 3.8.3, its
platform overhead, mediation and verification terms being the Class B,
Class C and Class D costs of this chapter. Where the declared loop
budget cannot absorb them, admission shall reject the agent or assign a
lower effective class as its class policy declares. An agent shall never
be admitted with a guardrail omitted.}
\reqfield{Dependency}{%
Req-T23-FW-20}
\reqfield{Traceability (Forward)}{%
\begin{itemize}
\tightlist
\item
  Arithmetic performed at admission by the DMO.
\end{itemize}}
\reqfield{UNITY-6G components}{%
\begin{itemize}
\tightlist
\item
  Admission (IDMO/DMO)
\item
  DMO
\end{itemize}}
\reqfield{Linked KPIs}{%
\begin{itemize}
\tightlist
\item
  KPI\#34: Reduce Open RAN real-time control loops time scales below 10
  ms and 10x in NTN for real-time 6G.
\item
  KPI\#4: Real-time reconfiguration of IEEE 802.11 network domain in
  range smaller than 10 ms from taking the decision.
\end{itemize}}
\end{requirement}

\begin{requirement}{Req-T23-KP-03}{MAY}{Envelope tiering for hard-real-time loops}
\reqfield{Description / rationale}{%
For hard-real-time loops the fabric may issue a verdict over an action
envelope, being a bounded region of the action space with a validity
window and a rate limit, within which the DE actuates without a
per-command mediation hop. Actuation outside the envelope, or after
expiry, shall require a fresh synchronous verdict. Envelope grants and
their expiry shall be held in the registry, and the last mile shall
re-validate against the live envelope and never a cached copy.}
\reqfield{Dependency}{%
Req-T23-CG-06}
\reqfield{Traceability (Forward)}{%
\begin{itemize}
\tightlist
\item
  Envelope held in the registry.
\item
  Live envelope checked at the last mile.
\end{itemize}}
\reqfield{UNITY-6G components}{%
\begin{itemize}
\tightlist
\item
  Conflict fabric
\item
  DMO
\item
  ACT runtime
\end{itemize}}
\reqfield{Linked KPIs}{%
\begin{itemize}
\tightlist
\item
  KPI\#4: Real-time reconfiguration of IEEE 802.11 network domain in
  range smaller than 10 ms from taking the decision.
\item
  KPI\#34: Reduce Open RAN real-time control loops time scales below 10
  ms and 10x in NTN for real-time 6G.
\end{itemize}}
\end{requirement}

\begin{requirement}{Req-T23-KP-04}{MAY}{Verification tiering for jitter-bounded loops}
\reqfield{Description / rationale}{%
Verification may be deferred for an action type whose action space
declares deferred timing and a verification deadline. The command shall
then proceed and verification complete within that deadline, which the
ACT runtime shall enforce as a rule granted with the admitted
descriptor, in the same manner as an action envelope. The action shall
report acknowledged\_\allowbreak{}unverified until confirmed, and a
missed deadline shall be a conformance fault feeding DE re-decision and
assurance. Tiering shall change when verification occurs and never
whether it occurs.}
\reqfield{Dependency}{%
Req-T23-ACT-04}
\reqfield{Traceability (Forward)}{%
\begin{itemize}
\tightlist
\item
  Verification deadline declared and enforced by the runtime.
\end{itemize}}
\reqfield{UNITY-6G components}{%
\begin{itemize}
\tightlist
\item
  ACT runtime
\item
  DE runtime
\item
  Evaluation
\end{itemize}}
\reqfield{Linked KPIs}{%
\begin{itemize}
\tightlist
\item
  KPI\#8: Cross-domain jitter of max of 2 milliseconds, with less than
  10\% of the packets missing deadline.
\item
  KPI\#4: Real-time reconfiguration of IEEE 802.11 network domain in
  range smaller than 10 ms from taking the decision.
\end{itemize}}
\end{requirement}

\begin{requirement}{Req-T23-KP-05}{MUST}{Recurring cost stays near-constant}
\reqfield{Description / rationale}{%
Mediation shall be per domain and not global. Per-service and per-agent
topics shall scale with services and agents, and agents shall not
cross-subscribe. Registry reads on the intent-egress path shall be
cache-backed with generation-fenced invalidation. Admission cost shall
be paid once at instantiation and not per cycle.}
\reqfield{Dependency}{%
Req-T23-IC-07}
\reqfield{Traceability (Forward)}{%
\begin{itemize}
\tightlist
\item
  Enforced by the placement of checks at admission.
\item
  Cache invalidation fenced by generation.
\end{itemize}}
\reqfield{UNITY-6G components}{%
\begin{itemize}
\tightlist
\item
  DMO
\item
  SMB
\item
  Platform runtime
\end{itemize}}
\reqfield{Linked KPIs}{%
\begin{itemize}
\tightlist
\item
  KPI\#37: Guarantee 99.99999\% service continuity through AI-driven
  proactive management of cloud-native services.
\item
  KPI\#6: Decrease service creation and termination time.
\item
  KPI\#24: Decrease end-to-end message delivery latency and increase
  maximum message publish throughput in the SBA service bus in
  integrated network scenarios.
\item
  KPI\#21: Reduce the performance penalty due to conflicts by 20\% by
  means of DLT, hierarchical algorithms and policy-based conflict
  resolution schemes.
\end{itemize}}
\end{requirement}

\begin{requirement}{Req-T23-KP-06}{MUST}{Fail closed stops deciding, never stops carrying traffic}
\reqfield{Description / rationale}{%
Loss of the mediation binding shall suspend decision emission and shall
not suspend the service; existing configuration shall persist and the
data plane shall continue. The agent shall move to suspended rather than
terminate. An action envelope valid at the loss shall remain valid until
its declared expiry, and an agent holding one shall be suspended at that
expiry if the binding has not been restored.}
\reqfield{Dependency}{%
Req-T23-CG-04}
\reqfield{Traceability (Forward)}{%
\begin{itemize}
\tightlist
\item
  Suspension applied by the DMO on binding loss.
\end{itemize}}
\reqfield{UNITY-6G components}{%
\begin{itemize}
\tightlist
\item
  DMO
\item
  Platform runtime
\end{itemize}}
\reqfield{Linked KPIs}{%
\begin{itemize}
\tightlist
\item
  KPI\#37: Guarantee 99.99999\% service continuity through AI-driven
  proactive management of cloud-native services.
\item
  KPI\#31: Reduce service recovery time below 180 s.
\end{itemize}}
\end{requirement}

\begin{requirement}{Req-T23-KP-07}{MUST}{Assurance preserves tail and variance guarantees}
\reqfield{Description / rationale}{%
Assurance returned to the requester shall not collapse to average
guarantees. Deadline-miss rate and jitter bound shall be reported as
distinct from mean latency, so that the report is usable against
deterministic and tail service level agreements.}
\reqfield{Dependency}{%
Req-T23-LC-08}
\reqfield{Traceability (Forward)}{%
\begin{itemize}
\tightlist
\item
  Reported per guarantee type by the assurance path.
\end{itemize}}
\reqfield{UNITY-6G components}{%
\begin{itemize}
\tightlist
\item
  DMO
\item
  IDMO
\item
  Intent layer
\end{itemize}}
\reqfield{Linked KPIs}{%
\begin{itemize}
\tightlist
\item
  KPI\#36: Guarantee 99.99999\% service availability and reliability
  based on effective monitoring and LCM.
\item
  KPI\#8: Cross-domain jitter of max of 2 milliseconds, with less than
  10\% of the packets missing deadline.
\end{itemize}}
\end{requirement}

\hypertarget{contribution-to-the-project-kpis-8}{%
\subsection{CONTRIBUTION TO THE PROJECT KPIs}\label{contribution-to-the-project-kpis-8}}

Table 12-4 states what the cost accounting of this chapter contributes
to each KPI it bears on, and Table 12-5 maps the requirements of this
chapter onto those KPIs.

\begin{wptable}{12-4}

\begin{longtable}[]{@{}
  >{\raggedright\arraybackslash}p{(\columnwidth - 4\tabcolsep) * \real{0.1499}}
  >{\raggedright\arraybackslash}p{(\columnwidth - 4\tabcolsep) * \real{0.6375}}
  >{\raggedright\arraybackslash}p{(\columnwidth - 4\tabcolsep) * \real{0.2126}}@{}}
\caption{Contribution of platform cost accounting to the project
KPIs}\tabularnewline
\toprule\noalign{}
\begin{minipage}[b]{\linewidth}\raggedright
\textbf{KPI}
\end{minipage} & \begin{minipage}[b]{\linewidth}\raggedright
\textbf{Contribution of cost accounting}
\end{minipage} & \begin{minipage}[b]{\linewidth}\raggedright
\textbf{Requirements}
\end{minipage} \\
\midrule\noalign{}
\endfirsthead
\toprule\noalign{}
\begin{minipage}[b]{\linewidth}\raggedright
\textbf{KPI}
\end{minipage} & \begin{minipage}[b]{\linewidth}\raggedright
\textbf{Contribution of cost accounting}
\end{minipage} & \begin{minipage}[b]{\linewidth}\raggedright
\textbf{Requirements}
\end{minipage} \\
\midrule\noalign{}
\endhead
\bottomrule\noalign{}
\endlastfoot
\textbf{KPI\#34} & Cost classification is the mechanism by which the
sub-10 ms loop is met: inference on the per-decision path, every other
guardrail off it. & Req-T23-KP-01, Req-T23-KP-02 \\
\textbf{KPI\#4} & Envelope tiering separates the authorizing decision
from the executing action, so a hard-real-time loop actuates at line
rate inside authority already granted. & Req-T23-KP-03 \\
\textbf{KPI\#8} & Verification tiering keeps read-back off the
jitter-bounded path, deferring when the system learns rather than
whether it checks. & Req-T23-KP-04 \\
\textbf{KPI\#37} & Fail closed means stop deciding, never stop carrying
traffic: agents are suspended rather than terminated, so recovery
restores an instance rather than re-instantiating it, and a valid action
envelope carries a hard-real-time loop through a brief fabric outage. &
Req-T23-KP-06 \\
\textbf{KPI\#21} & Per-domain mediation keeps conflict resolution
horizontally scalable as domains multiply. & Req-T23-KP-05 \\
\textbf{KPI\#6} & Admission cost paid once at instantiation rather than
per cycle protects service creation time from guardrail overhead. &
Req-T23-KP-05 \\
\textbf{KPI\#36} & Assurance preserves tail and variance guarantees
rather than collapsing to averages, so availability and reliability
claims can be checked against a deterministic or tail service level
agreement. & Req-T23-KP-07 \\
\end{longtable}

\end{wptable}

\begin{wptable}{12-5}\footnotesize

\begin{longtable}[]{@{}
  >{\raggedright\arraybackslash}p{(\columnwidth - 18\tabcolsep) * \real{0.1813}}
  >{\raggedright\arraybackslash}p{(\columnwidth - 18\tabcolsep) * \real{0.0910}}
  >{\raggedright\arraybackslash}p{(\columnwidth - 18\tabcolsep) * \real{0.0910}}
  >{\raggedright\arraybackslash}p{(\columnwidth - 18\tabcolsep) * \real{0.0910}}
  >{\raggedright\arraybackslash}p{(\columnwidth - 18\tabcolsep) * \real{0.0910}}
  >{\raggedright\arraybackslash}p{(\columnwidth - 18\tabcolsep) * \real{0.0910}}
  >{\raggedright\arraybackslash}p{(\columnwidth - 18\tabcolsep) * \real{0.0910}}
  >{\raggedright\arraybackslash}p{(\columnwidth - 18\tabcolsep) * \real{0.0910}}
  >{\raggedright\arraybackslash}p{(\columnwidth - 18\tabcolsep) * \real{0.0910}}
  >{\raggedright\arraybackslash}p{(\columnwidth - 18\tabcolsep) * \real{0.0911}}@{}}
\caption{Mapping of cost accounting requirements onto the project
KPIs}\tabularnewline
\toprule\noalign{}
\begin{minipage}[b]{\linewidth}\raggedright
\textbf{KPI\textbackslash{}}

\textbf{Requirement}
\end{minipage} & \begin{minipage}[b]{\linewidth}\raggedright
\textbf{\#4}
\end{minipage} & \begin{minipage}[b]{\linewidth}\raggedright
\textbf{\#6}
\end{minipage} & \begin{minipage}[b]{\linewidth}\raggedright
\textbf{\#8}
\end{minipage} & \begin{minipage}[b]{\linewidth}\raggedright
\textbf{\#21}
\end{minipage} & \begin{minipage}[b]{\linewidth}\raggedright
\textbf{\#24}
\end{minipage} & \begin{minipage}[b]{\linewidth}\raggedright
\textbf{\#31}
\end{minipage} & \begin{minipage}[b]{\linewidth}\raggedright
\textbf{\#34}
\end{minipage} & \begin{minipage}[b]{\linewidth}\raggedright
\textbf{\#36}
\end{minipage} & \begin{minipage}[b]{\linewidth}\raggedright
\textbf{\#37}
\end{minipage} \\
\midrule\noalign{}
\endfirsthead
\toprule\noalign{}
\begin{minipage}[b]{\linewidth}\raggedright
\textbf{KPI\textbackslash{}}

\textbf{Requirement}
\end{minipage} & \begin{minipage}[b]{\linewidth}\raggedright
\textbf{\#4}
\end{minipage} & \begin{minipage}[b]{\linewidth}\raggedright
\textbf{\#6}
\end{minipage} & \begin{minipage}[b]{\linewidth}\raggedright
\textbf{\#8}
\end{minipage} & \begin{minipage}[b]{\linewidth}\raggedright
\textbf{\#21}
\end{minipage} & \begin{minipage}[b]{\linewidth}\raggedright
\textbf{\#24}
\end{minipage} & \begin{minipage}[b]{\linewidth}\raggedright
\textbf{\#31}
\end{minipage} & \begin{minipage}[b]{\linewidth}\raggedright
\textbf{\#34}
\end{minipage} & \begin{minipage}[b]{\linewidth}\raggedright
\textbf{\#36}
\end{minipage} & \begin{minipage}[b]{\linewidth}\raggedright
\textbf{\#37}
\end{minipage} \\
\midrule\noalign{}
\endhead
\bottomrule\noalign{}
\endlastfoot
\textbf{Req-T23-KP-01} & \kpimark{} & & & & & & \kpimark{} & & \\
\textbf{Req-T23-KP-02} & \kpimark{} & & & & & & \kpimark{} & & \\
\textbf{Req-T23-KP-03} & \kpimark{} & & & & & & \kpimark{} & & \\
\textbf{Req-T23-KP-04} & \kpimark{} & & \kpimark{} & & & & & & \\
\textbf{Req-T23-KP-05} & & \kpimark{} & & \kpimark{} & \kpimark{} & & &
& \kpimark{} \\
\textbf{Req-T23-KP-06} & & & & & & \kpimark{} & & & \kpimark{} \\
\textbf{Req-T23-KP-07} & & & \kpimark{} & & & & & \kpimark{} & \\
\end{longtable}

\end{wptable}
